% 04d-noise.tex — v3 §3.5: the noise, per scale.
\subsection{The noise, per scale}\label{sec:noise-perscale}
% v3: §3.5

\paragraph{Inputs from \S\ref{sec:queue}.} Each is used frozen at the scale where it is read.
\begin{itemize}
\item The queue bound $Q\le B(u)$ (Theorem~\ref{thm:queue}, \eqref{eq:queueB}); so $|\Delta Q(W)|\le B(u)$ on windows at
  scale $u$.
\item The per-prime increment bound (Proposition~\ref{prop:ecl}, Corollary~\ref{cor:increment}). On $E_{\mathrm{cl}}^{(u)}$,
  for every slot window $I\subset[u,2u]$ with $|I|\le H(u)$, every non-absorbed latest prime $q\ge y(u)$, every $a\ge1$ and
  every jitter offset, $X_q^{(a)}(I)\le2M_1(u)$. The event $E_{\mathrm{cl}}^{(u)}$ fails on $\Good$ with probability
  $\le L_u^{-2}$.
\item The local count (Corollary~\ref{cor:localcount}). For dyadic $P$ and windows $J\subset[P,4P]$ of length $\le6H(P)$
  (part (i)), the non-absorbed count, and so the prime count, is $\le13.4|J|+K_{\mathrm{loc}}(P)+2$, with
  $K_{\mathrm{loc}}(P)=2B(4P)+7\lambda_P/3+23$ and $\lambda_P=4\log P$. For real $P\ge X_0$ (part (ii)) the same holds with
  $K_{\mathrm{loc}}(P)+K_{\mathrm{loc}}(2P)+4$ in place of $K_{\mathrm{loc}}(P)+2$.
\item Absorption (Lemmas~\ref{lem:lemmaP}, \ref{lem:rhoA}, Proposition~\ref{prop:S1S2}). Non-absorbed latest primes satisfy
  $q\ge y(u)=H(u)L_0L_u^6$, so $|I/q|\le L_u^{-6}$ for $|I|\le H(u)$. The small-$q$ case S1 does not arise in the present
  regime.
\end{itemize}

\paragraph{The composites with $q^2\mid n$: the source $s_2$.} Recall $C^{(\ge2)}=\Comp^{(2)}+\Mart^{(2)}$
(\S\ref{sec:noise-identity}). Fix $u$ and a window $I\subset[u,2u]$, and put $\sigma:=\frac12\sqrt{L_u\log L_u}$.
% Logged correction 52 (CORRECTIONS-vs-v3.md): v3 bounds the Dickman factor by 1/Gamma(max(1, L_u/s - 2)), which exceeds 1 on
% (1,2) (up to 1.129), and starts its sup at log X_0, although latest primes can lie in [X_0^-, X_0); it also applies the
% smooth-element weight with a smaller Dickman variable and without its factor e^{a_V}. The bound is re-derived below, with
% the endpoint X_0^-, and s_2 is defined by its closed form.
\begin{itemize}
\item The latest-prime filter forces $\tau_m<t(q)$ for $q^am$, so $\Pplus(m)\le2q$.
\item \emph{Interchange the sums.} Fix $a\ge2$ and $m$. Then $q^am\in I$ means $q\in I_{m,a}:=(I/m)^{1/a}$, and the $q$-sum of
  the landing probabilities is $\int_{I_{m,a}}\rhot\le1.01\,|I_{m,a}\cap[\Xm,\infty)|$, since $\rhot\le1.01$ surely (see the
  variance bound below) and $\rhot=0$ below $\Xm$. With $y=q^am\in I$ we have $dq=dy/(amq^{a-1})$, and on
  $I_{m,a}\cap[\Xm,\infty)$, $q\ge Q_{m,a}:=\max(\Xm,(u/m)^{1/a})$. Hence
  \[
    \frac{\Comp^{(2)}(I)}{|I|}\le1.01\sum_{a\ge2}\frac1a\sum_m\frac1{m\,Q_{m,a}^{a-1}},
  \]
  where $m$ runs over $\NN$ with $m\le2u/(\Xm)^a$ and $\Pplus(m)\le2(2u/m)^{1/a}$.
\item \emph{Blocks.} Split the $m$-sum into blocks $m\in(Y,2Y]$, $Y=2^{k-1}$ ($k\ge0$); for each $a$, at most $1.443L_u$ of
  them are nonempty. Define $s$ by $e^{as}=u/(2Y)$ and put $s_Y:=\max(\log\Xm,s)$. In the block, $Q_{m,a}\ge e^{s_Y}$,
  $\log Y=L_u-as-\log2$, and the $m$ are $z$-smooth with $\log z:=s+\log2+(2\log2)/a\le s+1.39$. A nonempty block has
  $Y<2u/(\Xm)^a$, so $s>\log\Xm-0.7$ and $\log z\ge\log\Xm-0.01$. The block contributes at most
  $\frac1ae^{-(a-1)s_Y}v(\log Y,z)$, where $v(\ell,z)$ is the $1/n$-weight of the $z$-smooth $n\in\NN\cap(e^\ell,2e^\ell]$, as in
  Lemma~\ref{lem:smooth-v}.
\item \emph{First bound.} Always $v\le1$ (the small heights in the proof of Lemma~\ref{lem:smooth-v}; the block $Y=\frac12$
  holds only $m=1$), and $s_Y\ge\log\Xm$. So
  $\Comp^{(2)}(I)/|I|\le1.01\cdot1.443L_u(\Xm)^{-1}\bigl(\frac12+\sum_{a\ge3}(\Xm)^{2-a}/a\bigr)\le0.73L_u/\Xm$.
\item \emph{Second bound, for $\sigma>\log\Xm$.} Then $L_u\ge10^{15}$, since $\log\Xm\ge10^8-1$. We show that every block term
  is at most $\frac1ae^{-\sigma}(\Xm)^{-(a-2)/2}$, that is, $(a-1)s_Y-\log v\ge\sigma+\frac12(a-2)\log\Xm$.
  \begin{itemize}
  \item If $s_Y\ge\sigma$, or if $(a-2)s_Y\ge2\sigma$, this holds with $v\le1$.
  \item Otherwise $s\le s_Y<\sigma$ and $as<4\sigma$, so $\ell:=\log Y\ge L_u-4\sigma-0.7\ge(1-4\cdot10^{-7})L_u$ and
    $\log z<\sigma+1.39$. Since $(a-1)s_Y\ge\frac12(a-2)\log\Xm$, it suffices that $-\log v\ge\sigma$. Put $u':=\ell/\log z$
    and $\omega:=2\sqrt{L_u/\log L_u}$, so that $\sigma=\frac14\omega\log L_u$ and
    $u'\ge(L_u-4\sigma-0.7)/(\sigma+1.39)\ge(1-4\cdot10^{-7})\,\omega$. For every $w\ge1$, iterating $\rho(w)\le\rho(w-1)/w$
    (from \ref{R1}) down to $w-\lfloor w\rfloor\in[0,1)$, where $\rho=1$, gives
    $\rho(w)\le\Gamma(w-\lfloor w\rfloor+1)/\Gamma(w+1)\le1/\Gamma(w+1)\le e^{-w\log w+w}$ (Stirling).
  \item If $u'\le\log z$: Lemma~\ref{lem:smooth-v} gives $v\le C_0\rho(u')e^{a_V}$ with
    $a_V=12.8\,u'\log(u'+2)/\log z\le1.3\cdot10^{-7}u'\log(u'+2)$. So
    \[
      -\log v\ge u'(\log u'-1)-1.3\cdot10^{-7}u'\log(u'+2)\ge(1-10^{-6})\,\omega\bigl(\tfrac12\log L_u-\tfrac12\log\log L_u-0.31\bigr),
    \]
    which is $\ge1.75\sigma$ for $\log L_u\ge34.5$: with $x:=\log L_u$ the ratio to $\sigma$ is $(1-10^{-6})(2-(2\log x+1.24)/x)$,
    which increases in $x$, and $\sigma\ge\log X_0^-$ forces $L_u\ge10^{15}$.
  \item If $u'>\log z$: then $\sqrt\ell>\log z$, and $v(\ell,z)\le v(\ell,z')$ with $\log z':=\sqrt\ell$. Lemma~\ref{lem:smooth-v} at
    the Dickman variable $\sqrt\ell$, where $a_V=12.8\log(\sqrt\ell+2)$, gives
    $-\log v\ge\sqrt\ell(\log\sqrt\ell-1)-12.8\log(\sqrt\ell+2)\ge5.5\sigma$; its ratio to $\sigma$ also increases in $\log L_u$, and
    is $\ge5.53$ at $L_u=10^{15}$.
  \end{itemize}
  Summing, $\Comp^{(2)}(I)/|I|\le1.01\cdot1.443L_u\,e^{-\sigma}\bigl(\frac12+\sum_{a\ge3}(\Xm)^{-(a-2)/2}/a\bigr)\le0.73L_u\,e^{-\sigma}$.
\item Together, for every window $I\subset[u,2u]$,
  \begin{equation}\label{eq:s2}
    \frac{\Comp^{(2)}(I)}{|I|}\le s_2(u):=9L_u\exp\Big(-\max\big(\log\Xm,\tfrac12\sqrt{L_u\log L_u}\big)\Big).
  \end{equation}
  In particular $s_2(u)\le9L_u/\Xm\le9(1+10^{-9})L_u/X_0$ always, since $X_0/\Xm\le1+H(X_0)/X_0$.
\item $\sum_{u\text{ dyadic}}\log(2u)\,s_2(u)\le e^{-9\cdot10^7}$, and the same holds with the weight $(\log3u)^2$. For
  $L_u\le\Lambda^2$, use $s_2(u)\le9L_u/\Xm$: these at most $1.443\Lambda^2$ scales contribute $\le14\Lambda^8/\Xm$. For
  $L_u>\Lambda^2$, $\frac12\sqrt{L_u\log L_u}\ge3.03\sqrt{L_u}$, and $10L^3e^{-3.03\sqrt L}$ is decreasing, so these scales
  contribute $\le10\Lambda^6e^{-3.03\Lambda}+\frac{10}{\log2}\int_{\Lambda^2}^\infty L^3e^{-3.03\sqrt L}\,dL\le10\Lambda^7e^{-3.03\Lambda}$.
  The total is $\le e^{-0.99\Lambda}$.
\item Put $A_u:=\#\{a\ge2:(\Xm)^a\le2u\}\le1.01L_u/\Lambda$. By Corollary~\ref{cor:increment} (which covers every
  $a\ge1$), the increment of a prime $q$ to $\Mart^{(2)}(I)$ is at most $b_2:=2A_uM_1(u)$.
\end{itemize}

\paragraph{The per-scale Freedman events.} Put $\lambda'_u:=3\log u+2\log\log u+2+\log H(u)$. For dyadic $u\ge X_0$, let
$\mathcal W'_u$ be the set of windows $[a,b]\subset[u,2u]$ with integer endpoints and $b-a\le1.03H(u)+1$; there are at most
$(u+1)(1.03H(u)+2)$ of them. Let $\mathcal W_u\subset\mathcal W'_u$ consist of those with $b-a\ge H(u)-1$.
\begin{itemize}
\item $G_u^{\mathrm{NV}}$ (the event E3 of Definition~\ref{def:good}): every $W\in\mathcal W_u$ satisfies \eqref{eq:F1}
  and \eqref{eq:F2}.
\item $G_u^{\mathrm{sw}}$ (the event E3$'$): every $W\in\mathcal W'_u$ satisfies \eqref{eq:F1} and \eqref{eq:F2}.
\end{itemize}
\begin{align}
  |\Mart(W)|&\le\sqrt{2\bar V_W\lambda'_u}+\tfrac43M_1(u)\lambda'_u,&\bar V_W&:=2M_1(u)\cdot50\log(2u)\cdot|W|;\label{eq:F1}\\
  |\Mart^{(2)}(W)|&\le\sqrt{2\bar V^2_W\lambda'_u}+\tfrac23b_2\lambda'_u,&\bar V^2_W&:=b_2\cdot s_2(u)\cdot|W|.\label{eq:F2}
\end{align}

\paragraph{Truncation.} As in Theorem~\ref{thm:queue}(i), Freedman's inequality is applied to the truncated increments
$\tilde X_q(W):=\min(X_q(W),2M_1(u))$ and $\tilde X^{(2)}_q(W):=\min(X^{(2)}_q(W),b_2)$, each compensated by its own
conditional mean. These are martingales in decision order, with increments $\le2M_1$ and $\le b_2$ surely. Freedman's
inequality is used in the form
\[
  \Prob(\exists n:\ S_n\ge a,\ V_n\le v)\le\exp\Big(-\frac{a^2}{2(v+ba/3)}\Big),
\]
for increments bounded by $b$ \cite[(4.1)]{Freedman1975} (with Bennett's bound for the exponent), with the deterministic $a$
and $v=\bar V_W,\bar V^2_W$ of \eqref{eq:F1}, \eqref{eq:F2}.
\begin{itemize}
\item \emph{Truncation is inactive on $\Good(<t(W))$.} For every latest prime $q$ contributing to $W$, the set of cofactors
  (elements whose primes are all decided before $q$) is $\mathcal F_{q-}$-measurable. So ``$X_q(W)\le2M_1$ for every jitter
  offset of $q$'' is an $\mathcal F_{q-}$-measurable event. It may not yet be certified at time $s_q$ (for $q<\sqrt{2u}$).
  But on $\Good(<t(W))$ the event $E_{\mathrm{cl}}^{(u)}$ holds (its time $8u/L^{11}+H$ is $<t(W)$), and
  Corollary~\ref{cor:increment} then gives it for every $q$ and every offset. So $\tilde X_q=X_q$ and
  $\E[\tilde X_q\mid\mathcal F_{q-}]=\E[X_q\mid\mathcal F_{q-}]$. The truncated and the true martingales therefore agree on
  $\Good(<t(W))$, and likewise for $\Mart^{(2)}$.
\item \emph{The variance bound holds on $\Good(<t(W))$.} We have
  $\mathrm{Var}(\tilde X_q\mid\mathcal F_{q-})\le2M_1\E[X_q\mid\mathcal F_{q-}]$, so $V_W\le2M_1\Comp(W)$; likewise
  $V^2_W\le b_2\Comp^{(2)}(W)\le b_2s_2|W|$. Next,
  $\Comp(W)=\int_W(\Kpt\rhot-K_{\Abs}\rhot+\NT_{\mathrm{pt}})\le1.01\sum_v\omega_y(v)\cdot|W|+\int_W|\NT_{\mathrm{pt}}|\le50\log(2u)\,|W|$.
% Logged correction 73: v3 says (H_BS) "at heights <= u"; a short window W may close at t(W) < u + 12 log u + 3, where that is
% not certified. Lemma S(i) at y in W reads only cofactors w <= y/v <= 4u/X_0, far below the horizon of Good(<t(W)).
  This uses $\sum_v\omega_y(v)\le1.2\log_2(4u/X_0)$ (Lemma~\ref{lem:S}(i), that is, \textup{(H$_{\mathrm{BS}}$)} at heights
  $\le4u/X_0$, below the certified horizon of $\Good(<t(W))$, Remark~\ref{rem:horizon}), and the sure bound $\rhot\le1.01$. The latter holds because each continuous
  slot decides at most one prime, and the starts of consecutive continuous slots are at least $1-10^{-8}$ apart (between them
  $\int T'_{\mathrm{cont}}=1$, \S\ref{sec:greedy}, with $T'_{\mathrm{cont}}\le1+10^{-8}$). No upper bound on the length of a slot is used
  in this section.
% Logged correction 77: one unit of continuous mass lies between consecutive continuous-slot starts, not in each continuous slot
% (a continuous slot followed by a jump slot, or an empty one in the tie case, carries less).
\item So on $\Good(<t(W))$ the event $\{|\Mart(W)|>$ the bound of \eqref{eq:F1}$\}$ lies in Freedman's event for the
  truncated martingale; likewise for \eqref{eq:F2}. No positivity of the margin is used.
\end{itemize}

\paragraph{Probability.} With $a=\sqrt{2v\lambda'_u}+\frac23b\lambda'_u$, Freedman's inequality gives, per window and per
martingale, a two-sided failure probability $\le2e^{-\lambda'_u}=2e^{-2}u^{-3}L_u^{-2}H(u)^{-1}$. With four bounds per window,
\[
  \Prob\big((G_u^{\mathrm{NV}}\cap G_u^{\mathrm{sw}})^c\cap\Good(<t)\big)\le8e^{-2}(u+1)(1.03H(u)+2)u^{-3}L_u^{-2}H(u)^{-1}
  \le1.2\,u^{-2}L_u^{-2}.
\]
Each window's bound is an event whose time is the moment the window closes ($\le2u$). These are the events E3 and E3$'$ of
Definition~\ref{def:good}. Together with $E_{\mathrm{cl}}$ they enter $\Good(<t)$ and the budget $\Prob(\tau<\infty)\le1.5/\Lambda$
of Lemma~\ref{lem:budget}. The budget is dominated by $E_{\mathrm{cl}}$: $\sum_uL_u^{-2}\le1.45/\Lambda$, while the
Freedman part is $\le2X_0^{-2}$.

\begin{lemma}[per-scale noise]\label{lem:noise}
% v3: Lemma 3.11
For $y\in[u,2u]$, work on $\Good(<y+H(y))$. This event contains the E3/E3$'$ windows of the blocks
$u'\in\{u/2,u,2u\}$ that close by $y+H(y)/2$, and $E_{\mathrm{cl}}^{(u')}$. Put
\begin{multline*}
  \nu_0(u'):=\sqrt{\frac{206M_1(u')\lambda'_{u'}\log(2u')}{H(u')}}\Big(1+\frac1{H(u')}\Big)
  +\sqrt{\frac{4.2A_{u'}M_1(u')s_2(u')\lambda'_{u'}}{H(u')}}\\
  +\frac{\frac43(1+A_{u'})M_1(u')\lambda'_{u'}+3B(u')+100\log(2u')+6}{H(u')}+1.1s_2(u').
\end{multline*}
Then
\[
  |\nu(y)|\le\bar\nu(u):=2.05\max_{u'\in\{u/2,u,2u\}}\nu_0(u').
\]
% Logged correction 73 (R17, C73-6): the good event of the bound on E(I) is stated.
For every slot window $I\subset[u,2u]$ with $|I|\le H(u)$, put $E(I):=\sigma(I)-\int_Im$, where
$m=T'-\Kpt\rhot-\NT_{\mathrm{pt}}$ is the pointwise margin density (\S\ref{sec:noise-identity}). Then, on $\Good(<\max I+H(u))$,
\begin{multline*}
  |E(I)|\le\sqrt{200M_1(u)\lambda'_u\log(2u)(|I|+2)}+\sqrt{4A_uM_1(u)s_2(u)\lambda'_u(|I|+2)}\\
  +\tfrac43(1+A_u)M_1(u)\lambda'_u+B(u)+100\log(2u)+s_2(u)(|I|+2)+2 .
\end{multline*}
Write $\hat\nu(u):=\sup_{|I|\le H(u)}|E(I)|/H(u)$. Both bounds depend on $u$ only: nothing at the current scale $x$
enters.
\end{lemma}

\begin{proof}
\emph{The bound on $\nu$.} By \eqref{eq:margin-identity},
$H(y)\nu(y)=\Mart(W_y)+\Mart^{(2)}(W_y)-\Delta Q(W_y)+\Comp^{(2)}(W_y)+E_{\mathrm{edge}}$.
\begin{itemize}
\item $E_{\mathrm{edge}}$ counts the two partial slot cells at the ends of $W_y$. Each holds at most one slot and at most
  $A_j\le B(u)+1$ composites, so $|E_{\mathrm{edge}}|\le2B(u)+4$.
\item $|\Delta Q|\le B(u)$ by Theorem~\ref{thm:queue}, and $\Comp^{(2)}(W_y)\le s_2|W_y|$ by \eqref{eq:s2}.
\item Choose $W^-\subset W_y\subset W^+$ in $\mathcal W_u$ with $|W^+\setminus W^-|\le2$. Counts and compensators are
  monotone in the window, so $\Mart(W^-)-\Comp(W^+\setminus W^-)\le\Mart(W_y)\le\Mart(W^+)+\Comp(W^+\setminus W^-)$, with
  $\Comp(W^+\setminus W^-)\le100\log(2u)$. The same holds for $\Mart^{(2)}$, with $2s_2$ in place of $100\log(2u)$.
\item Apply \eqref{eq:F1} and \eqref{eq:F2} with $|W^\pm|\le1.03H(u)+1$, and divide by $H(y)\ge H(u)$.
\item If $W_y$ crosses $u$ or $2u$ (that is, $y$ is within $H(y)/2$ of it), split $W_y$ there. Each piece lies in one
  dyadic block $u'$ and has length $\le1.03H(u')$, so the above applies to it in block $u'$. The two bounds add, and
  $H(y)/H(u')\le1.02$ by \eqref{eq:RH}; the factor $2.05$ covers this.
\end{itemize}

\emph{The bound on $E(I)$.} The bookkeeping \eqref{eq:bookkeeping} on the slot window $I$ has no partial cells. By the
computation of \eqref{eq:margin-identity} without the $B$-average,
$\sigma(I)=\int_Im-\Comp^{(2)}(I)-\Mart(I)-\Mart^{(2)}(I)+\Delta Q(I)\pm2$. Now sandwich $I$ between windows of
$\mathcal W'_u$ as above, and apply \eqref{eq:F1}, \eqref{eq:F2} and \eqref{eq:s2}.
\end{proof}

\paragraph{The \textup{(SW)} sums.} Under the tolerance $M_1\le H/L^{14}$ of \eqref{eq:params}, with $B\le28M_1L^2$,
$A_u\le1.01L_u/\Lambda$ and $\lambda'_u\le18.2L_u$ (since $\log H\le15.1\Lambda$ on the raised range),
\[
  \bar\nu(u)\le128L_u^{-6}+2.3\max_{u'}s_2(u'),\qquad\hat\nu(u)\le62L_u^{-6}+1.1s_2(u),
\]
and
\begin{equation}\label{eq:sw-sums}
  \sum_{u\text{ dyadic}\ge X_0}\log(2u)\,\bar\nu(u/2)\le10^{-30},\qquad
  \sum_v(\log3v)^2\,\bar\nu(v/2)\le10^{-22},\qquad
  \sum_v\log(3v)\,\hat\nu(v)\le10^{-30}.
\end{equation}
\begin{itemize}
\item The leading term is $\sqrt{206\cdot18.2}\,L\,(M_1/H)^{1/2}\le61.3L^{-6}$. The linear terms are $\le140L^{-12}$:
  $\frac43(1+A_u)M_1\lambda'\le50L^2M_1$ and $3B\le84L^2M_1$. Then $\sum_j128(\Lambda+j\log2)^{-5}\le46.2/\Lambda^4$ and
  $\sum_j128(\Lambda+j\log2)^{-4}\le61.6/\Lambda^3$.
\item On the standard range, where $H=H_{\mathrm{std}}$ and $M_1\le\frac12H^{1/2}$, the same sums are $\le e^{-600}$; the first
  term is then $\approx18L\,e^{-2(\log L)^2}$.
\item The sums converge whenever $H(v)/M_1(v)\ge(\log v)^{8+\delta}$. This is what fixes the tolerance $M_1\le H/L^{14}$.
\end{itemize}
